% Folder 136, batch 10 (archivists' pages 181–200).
% Pass: Opus 5.5 (claude-opus-5-5), 2026-10-03 — first pass, unchecked against the pages by a human.
%
% The even pages (182, 184, …, 200) are versos carrying a typescript unrelated
% to these notes; they are skipped and do not appear. The odd pages are his
% manuscript, written in the gaps of the same paper.
\documentclass[11pt,a4paper]{article}
\input{../preamble/grothendieck}

\folder{136}
\batch{10}
\pages{181}{200}
\foldertitle{Complexe de De Rham à puissance divisée [conférence de 1976 à l’IHÉS] : notes manuscrites (s.d.).}
\dating{[à partir de 1975-1976]}
\watermark{Édition de démonstration}

\begin{document}

\page{181}
\note{la page commence au milieu d'un raisonnement venu des pages précédentes (« Si donc… »).}
Si donc on veut trouver des op.\ coh.\ intéressantes \ill{} les opérations de
$\Omega$ \ill{} sur le corps, il faut \ill{} se placer en \emph{Cas 2}
(\uncertain{en caractéristique}) -- on peut espérer trouver alors une
description des \struck{\ill{}} \add{carrés} de Steenrod mod 2 (regarder ce
qu'il en est…).

Mais on peut essayer d'être plus astucieux, en appliquant p.\ ex.\ le foncteur
\[
  \tau_\omega : (\text{Mod gr de } S_h \Longrightarrow)\ M \longrightarrow M_\omega ,
\]
donc $X_* \mapsto \operatorname{Hom}(X_*, A_h^{*})$ devient un foncteur sur
(Semisimpl.) à valeurs dans les objets \struck{gradués} \add{algèbres} diff.\
de $M_\omega$ (qui est une cat.\ abélienne « \uncertain{de Artin} »
\uncertain{avec} $\otimes$ \uncertain{connu}), et on peut se proposer, laissant
tomber la structure graduée au profit de la structure d'anneau diff.\ sans
plus, de déterminer les types des anneaux et des foncteurs (d'abord pour $h$
fixé, puis sa variation avec $h$). Peut-être trouverons-nous des op.\
\uncertain{coh.}\ intéressantes \ill{} puiss.\ de Steenrod…

\page{183}
\marginal{Endomorphismes de la $S$-algèbre différentielle semisimpliciale $A_*$}
\[
  A_1 = k\{X_0, X_1, dX_0, dX_1\}/(dX_0 + dX_1)_{\mathrm{pd}}
      = k\{T\}\{X_0, dX_0\} \simeq k\{T\}\{X_1, dX_1\},
\]
\[
  X_0 + X_1 = T .
\]
\note{le signe $\simeq$ de la deuxième égalité est surchargé.}
\[
  \varphi_1(X_0^{(i)}) = \xi_i ,
\]
\[
  (1)\qquad \boxed{\xi_i\,\xi_j = c_{ij}\,\xi_{i+j}} \qquad (i, j \geqslant 1)
\]
\[
  \varphi_1(dX_0) = d\xi_1
\]
\marginal{au crayon : \ill{} et $d\xi_1$ commute avec $\xi_i$ ($i \geqslant 1$) en $A_1$ commutatif !}
\note{à côté, une ligne au crayon et une rature hachurée, illisibles.}
\[
  \varphi_1(dX_0^{(i)}) \overset{?}{=} d\varphi_1(X_0^{(i)}) ,
\]
c'est-à-dire $\varphi_1(X_0^{(i-1)}dX_0) = d\xi_i$, et
$\varphi_1(X_0^{(i-1)}dX_0) = \xi_{i-1}\,d\xi_1$, d'où
\[
  (2)\qquad \boxed{d\xi_i = \xi_{i-1}\,d\xi_1} \qquad (i \geqslant 2).
\]
\marginal{au crayon : on aura $\varphi_1(d(X_0^{(i)}dX_0)) \overset{?}{=} d\varphi_1(X_0^{(i)}dX_0)$, avec $d(\xi_i\,d\xi_1) = 0$ et $\varphi_1(0) = 0$.}

\begin{tikzcd}
A_1 \arrow[r, bend left=20, "\alpha"] \arrow[r, bend right=20, "\beta"'] & A_0 = k\{T\}
\end{tikzcd}

\[
  \alpha(X_0^{(i)}) = 0 \ \ \forall i \geqslant 1, \quad \alpha(dX_0) = 0 ;
  \qquad
  \beta(X_0^{(i)}) = T^{(i)} \ \ \forall i \geqslant 1, \quad \beta(dX_0) = 0 .
\]
\[
  \begin{cases}
    \alpha \circ \varphi_1 = \alpha \\
    \beta \circ \varphi_1 = \beta
  \end{cases}
\]
$\alpha\circ\varphi_1(X^{(i)}) = \alpha(\xi_i)$ et $\alpha(X^{(i)}) = 0$ ;
$\beta\circ\varphi_1(X^{(i)}) = \beta(\xi_i)$ et $\beta(X^{(i)}) = T^{(i)}$ :
\[
  (3\mathrm{a})\quad \boxed{\alpha(\xi_i) = 0}
  \qquad\qquad
  (3\mathrm{b})\quad \boxed{\beta(\xi_i) = T^{(i)}}
\]
\[
  \xi_i = \overbrace{\sum_{j \geqslant 0} a_{ij} X_0^{(j)}}^{A_i}
        + \Bigl(\overbrace{\sum_{j \geqslant 0} b_{ij} X_0^{(j)}}^{B_i}\Bigr) dX_0 ,
  \qquad a_{ij}, b_{ij} \in k\{T\} .
\]
\note{les accolades $A_i$, $B_i$ sont au crayon.}
(3a) s'écrit $a_{i,0} = 0$ pour tout $i \geqslant 1$, i.e.
\[
  (3'\mathrm{a})\qquad \xi_i = \sum_{j \geqslant 1} a_{ij} X_0^{(j)} + \cdots
  \qquad \forall i \geqslant 1 ,
\]
(3b) s'écrit
\[
  (3'\mathrm{b})\qquad \sum_{j \geqslant 1} a_{ij}\,T^{(j)} = T^{(i)} \qquad \forall i \geqslant 1 .
\]
\note{le développement de (3'a) est entouré et souligné trois fois ; avant (3'b), une somme biffée.}
\[
  d\xi_i = \Bigl(\sum_{j \geqslant 1} a_{ij} X_0^{(j-1)}\Bigr) dX_0 ,
  \qquad
  d\xi_1 = \Bigl(\sum_{j \geqslant 1} a_{1j} X_0^{(j-1)}\Bigr) dX_0 ,
\]
\[
  \xi_{i-1} = \sum_{j \geqslant 1} a_{i-1,j} X_0^{(j)}
            + \Bigl(\sum_{j \geqslant 0} b_{i-1,j} X_0^{(j)}\Bigr) dX_0 ,
\]
\[
  \xi_{i-1}\,d\xi_1
  = \Bigl(\sum_{k \geqslant 1} a_{1k} X_0^{(k-1)}\Bigr)
    \Bigl(\sum_{k \geqslant 1} a_{i-1,k} X_0^{(k)}\Bigr) dX_0 .
\]
\note{le calcul s'arrête là ; la page suivante repart sur une formulation générale.}

\page{185}
Finalement, on trouve

\textbf{Proposition.} Soit $S$ un anneau commutatif, $A = S\{X_0\}[dX_0]$,
\struck{on cherche les} $S$-algèbre diff.\ à puiss.\ div., augmentation
$\alpha : A \to S$ ($\alpha(X_0^{(i)}) = 0$ pour $i \geqslant 1$,
$\alpha(dX_0) = 0$) (N.B.\ on regarde $S$ comme anneau diff.\ avec $d_S = 0$,
i.e.\ une puiss.\ div.\ compatible avec sa structure \ill{}). On cherche les
endomorphismes $\varphi$ de $S$-algèbre différentielle augmentée (on dit aussi
p.d.\,!) \struck{\ill{}} de $A$ :
\[
  \varphi : A \to A \ \text{endom.\ de $S$-algèbre}
  \quad [\varphi(x+y) = \varphi(x) + \varphi(y),\ \varphi(xy) = \varphi(x)\varphi(y),\ \varphi(1) = 1],
\]
\[
  \varphi\, d = d\, \varphi , \qquad \alpha\,\varphi = \alpha .
\]
Alors $\varphi$ est donné par les
\[
  \varphi(X_0^{(i)}) = \xi_i \in A \quad (i \geqslant 1),
  \qquad \xi_i = A_i + B_i\,dX_0 , \quad A_i, B_i \in S\{X_0\} \ (i \geqslant 1)
\]
(on aura $\varphi(dX_0) = d\varphi(X_0) = d\xi_1 = dA_1 = A_1'\,dX_0$), avec les conditions
\[
  \begin{cases}
    A_i = A_1^{(i)} \\
    \alpha(A_1) = 0 , \ \text{i.e. } A_1 \in S\{X_0\}^{+} .
  \end{cases}
\]
Donc $\varphi$ est donné par
\[
  \boxed{\begin{array}{l} A_1 \in S\{X_0\}^{+} \\ B_i \in S\{X_0\} \ (i \geqslant 1) \end{array}}
\]
sans autre condition !

Pour que $\varphi$ commute aux p.d., il faut et il suffit que l'on ait
$\xi_i = \xi_1^{(i)}$ ($i \geqslant 2$), ou encore
\[
  B_i = A_1^{(i-1)} B_1 .
\]

\page{187}
Donc les $\varphi$ correspondent exactement aux
$\xi_1$ \struck{$A_1 + B_1\,dX_0$} $\in \operatorname{Ker}\alpha = A^{+}$
(i.e.\ on choisit $\xi_1 = A_1 + B_1\,dX_0$, $A_1, B_1 \in S\{X_0\}$),
aux couples $A_1 \in S\{X_0\}^{+}$, $B_1 \in S\{X_0\}$.

\textbf{Cor.} Si $\varphi$ respecte la graduation et la structure d'algèbre
différentielle augmentée, alors $\varphi$ respecte les p.d.\,!

On prend un anneau $S = k\{T\}$, et on regarde $S$ comme anneau \add{$k\{T\}^{+}$}
p.d.\ par $A = A_1$. Donc $\alpha$ est compatible avec les structures, y
compris les p.d. On définit
\[
  \beta : A_1 = S\{X_0, dX_0\} = k\{T\}\{X_0, dX_0\} \longrightarrow S ,
\]
\[
  A_1 \simeq k\{X_0, X_1, dX_0, dX_1\}/\bigl((X_0 + X_1 - T,\ dX_0 + dX_1)\bigr)_{\mathrm{pd}} ,
\]
par la condition que $\beta$ soit comp.\ avec les structures, et $\beta(X_0) = T$, donc
\[
  \begin{cases}
    \beta(X_0^{(i)}) = T^{(i)} \\
    \beta(dX_0) = 0
  \end{cases}
\]
(car $dT = 0$ \ill{} : on doit avoir \ill{} diff.\ sur $S$).
\marginal{N.B.\ $\alpha(X_0) = 0$, $\alpha(X_1) = T$ ; $\beta(X_0) = T$, $\beta(X_1) = 0$.}

Et on s'intéresse aux $\varphi$ \struck{compatibles} \add{\ill{}} avec
l'augmentation $\alpha$ et $\beta$. Cela signifie que l'on a
$\beta\varphi = \beta$, i.e.\
$\beta\varphi(X_0^{(i)}) = \beta(X_0^{(i)})$ pour tout $i \geqslant 1$, où
$\varphi(X_0^{(i)}) = A_1^{(i)} + B_i\,dX_0$, ou encore
\[
  \beta(A_1)^{(i)} = T^{(i)} \qquad \forall i \geqslant 1 ,
\]
ou encore

\page{189}
\note{suite directe de la page 187.}
\[
  \boxed{\beta(A_1) = T} \qquad \text{i.e. } A_1(T) = T .
\]
\struck{Donnons} Considérons
\[
  A_1 - X_0 = \varphi^0_0(X_0) - X_0 = -\bigl(\varphi^0_0(X_1) - X_1\bigr) = f \in A_1^0
\]
(où $\varphi_0(x)$ est la composante de $\varphi(x)$ suivant
\[
  S\{X_0\} \simeq k\{T\}\{X_0, X_1\}/(X_0 + X_1 - T)_{\mathrm{pd}}
\]
dans
\[
  A_1 = \underbrace{S\{X_0\}}_{A_1^0} + \underbrace{S\{X_0\}\,dX_0}_{A_1^1}
      \simeq S\{X_1\} + S\{X_1\}\,dX_1 \ ).
\]
[N.B.\ Comme $X_1 = T - X_0$, on a
$\varphi^0_0(X_1) = \varphi^0_0(T) - \varphi^0_0(X_0) = T - \varphi^0_0(X_0)$,
donc $\varphi^0_0(X_1) - X_1 = X_0 - \varphi^0_0(X_0)$.]
\note{l'indice inférieur de $\varphi^0_0$ est peu net.}

Alors la conjonction des conditions $A_1(X_0) = 0$ et $A_1(T) = T$ s'écrit
aussi
\[
  f \in (X_0)_{\mathrm{pd}} \cap (X_1)_{\mathrm{pd}} ,
\]
donc, si l'on regarde $A_1^0$ comme \struck{produit de} $k\{X_0, X_1\}$, cela
signifie qu'on doit avoir
\[
  \boxed{f = \sum_{i, j \geqslant 1} f_{ij}\, X_0^{(i)} X_1^{(j)} , \quad f_{ij} \in k}
\]
d'où un résultat :

\textbf{Prop.} Les endom.\ $\varphi$ de la $S$-algèbre diff.\
\[
  A_1 \simeq S\{X_0, X_1, dX_0, dX_1\}/\!/(X_0 + X_1 - T,\ dX_0 + dX_1)_{\mathrm{pd}} ,
\]
où $S = k\{T\}$ (avec sa puiss.\ div.), qui sont compatibles aux deux
augm.\ $\alpha, \beta : A_1 \to S$ ($\alpha(X_0) = 0$, $\alpha(X_1) = T$ ;
$\beta(X_0) = T$, $\beta(X_1) = 0$) sont donnés exactement par les systèmes
$f \in A_1^0 = k\{X_0, X_1\}$ avec
\[
  f \in \operatorname{Ker}\bigl(k\{X_0, X_1\} \to k\{X_0\}\bigr)
    \cap \operatorname{Ker}\bigl(k\{X_0, X_1\} \to k\{X_1\}\bigr) ,
\]
\[
  \text{i.e.}\quad \boxed{f = \sum_{i, j \geqslant 1} c_{ij}\, X_0^{(i)} X_1^{(j)}}
\]

\page{191}
\note{suite directe de la page 189.}
et les
\[
  \boxed{B_i \in A_1^0 = k\{X_0, X_1\} \qquad (i \geqslant 0)}
\]
par les formules
\[
  \begin{cases}
    \varphi(X_0^{(i)}) = (X_0 + f)^{(i)} + B_i\,dX_0 \\
    \varphi(dX_0) = 1 + df
  \end{cases}
\]
\note{« $1 + df$ » est bien ce qu'on lit ; $d(X_0 + f) = dX_0 + df$ serait attendu.}
(on sait que $f = \varphi^0(X_0) - X_0 = -(\varphi^0(X_1) - X_1)$ ;
$B_i\,dX_0 = \varphi^1(X_0^{(i)})$).

\emph{N.B.}\ Si on échange le rôle de $X_0$ et $X_1$, $f$ change de signe,
-- qu'en est-il \uncertain{des} $B_i$ ? Il faut calculer
\[
  \varphi^1(X_1^{(i)}) = \varphi^1\bigl((T - X_0)^{(i)}\bigr)
  = \varphi^1\Bigl(\sum_{0 \leqslant j \leqslant i} T^{(i-j)} X_0^{(j)} (-1)^j\Bigr)
  = \Bigl(\sum_{0 \leqslant j \leqslant i} T^{(i-j)} (-1)^j B_j\Bigr) dX_0
\]
\[
  = \Bigl(\underbrace{\sum_{0 \leqslant j \leqslant i} (-1)^{j+1} T^{(i-j)} B_j}_{C_i}\Bigr) dX_1 .
\]
\struck{En tout cas $\varphi^1(T)$ —}
\[
  \varphi^1\bigl((X_0 + X_1)^{(i)}\bigr) = \varphi^1(T^{(i)}) = 0 ,
  \qquad
  (X_0 + X_1)^{(i)} = \sum_{j+h=i} X_0^{(j)} X_1^{(h)} ,
\]
\[
  \varphi^1\bigl((X_0 + X_1)^{(i)}\bigr)
  = \text{comp.\ de degré 1 de }\sum_{j+h=i} \varphi(X_0^{(j)})\,\varphi(X_1^{(h)}) ,
\]
\[
  \varphi(X_0^{(j)}) = (X_0 + f)^{(j)} + B_j\,dX_0 ,
  \qquad
  \varphi(X_1^{(h)}) = (X_1 - f)^{(h)} + C_h\,dX_1 , \quad dX_1 = -dX_0 ,
\]
donc
\[
  \sum_{j+k=i} \Bigl[ -(X_0 + f)^{(j)} C_k + B_j (X_1 - f)^{(k)} \Bigr] = 0 .
\]
\note{la ligne médiane porte, sous l'accolade, une notation « $K^1$ » pour la composante de degré 1 ; rendue ici en clair.}

\page{193}
\textbf{Prop.} Soit $\varphi$ un endom.\ de
\struck{$A_1 \simeq k\{X_0, X_1, dX_0, dX_1\}/\ill{}$}
\[
  A_1 = k\{X_0, X_1\}[dX_0, dX_1]/(dX_0 + dX_1)\ \struck{\ill{}} ,
\]
et $\varphi$ \struck{\ill{}} un endom.\ de $S$-algèbre différentielle
($S = k\{T\}$) compatible aux deux augmentations $\alpha$ et $\beta$ (qui en font
des augm.\ simpliciales $A_1 \to A_0$\,!). Soit $n \geqslant 0$ ; je dis qu'il
existe un et un seul endom.\ $\varphi_n$ de $S$-alg.\ différentielle de $A_n$,
telle que pour toute application simpliciale (croissante\,!)
$\sigma : \Delta_n \to \Delta_1$, l'hom.\ $A_*(\sigma) : A_1 \to A_n$ satisfasse
\[
  \struck{\ill{}} \quad \varphi_n\, A_*(\sigma) = A_*(\sigma)\, \varphi_1
\]
\struck{si $n$} (si $n \geqslant 1$, il suffit \struck{\ill{}} \ill{} des
celles \ill{} aux $\sigma$ croissantes).

\emph{Cor.}\ \struck{Donc $\varphi_1 = \varphi$, et les syst.\
$\varphi_* = (\varphi_n)$} Il y a un endom.\ de la $S$-algèbre différentielle
semi-simpliciale $A_*$, caractérisé par le fait que $\varphi_1 = \varphi$.
L'ens.\ des endom.\ de $S$-alg.\ diff.\ $A_*$ est en correspondance
biunivoque (conservant les compositions) avec l'ens.\ des $\varphi$
(\uncertain{compatibles à la p.d.}).
\note{un trait vertical en marge gauche réunit l'énoncé du corollaire ; un trait oblique le barre en partie, renvoyant à la note marginale « faux tel quel ».}
\marginal{faux tel quel ; vrai si $2B_i = 0$ $\forall i \geqslant 1$, i.e.\ $2\varphi^1(X_0^{(i)}) = 0$ $\forall i \geqslant 1$ \ill{} $2\cdot 1_k = 0$ ; c'est-à-dire si $\varphi$ conserve la \ill{} $A_1^0$ qui soit comp.\ avec $\alpha^0, \beta^0$)}
\marginal{en bas, en oblique : \ill{}}

\emph{Dém.\ du corollaire.} (Il suffit de montrer que si
$\tau : \Delta_{n'} \to \Delta_n$, d'où $A_*(\tau) : A_n \to A_{n'}$, alors
\[
  \varphi_{n'}\, A_*(\tau) = A_*(\tau)\, \varphi_n .
\]
(C'est clair si $n = 0$, car $A_0 = S$, et par hyp.\ si $n = 1$.
\uncertain{On peut supposer} $n \geqslant 1$.)
Or il suffit de \ill{} en composant avec les $\sigma : \Delta_n \to \Delta_1$
\struck{\ill{}} \add{\ill{}} les $A_*(\sigma) : A_1 \to A_n$ (\ill{} \ill{}
\uncertain{engendrent} les algèbres $A_n$ dans \ill{} sens \ill{} \ill{}

\page{195}
\note{suite directe de la page 193.}
\[
  \varphi_{n'}\, \underbrace{A_*(\tau)\, A_*(\sigma)}_{A_*(\sigma\tau)}
  = A_*(\tau)\, \underbrace{\varphi_n\, A_*(\sigma)}_{A_*(\sigma)\,\varphi \ (\text{par hyp.\ sur } \varphi_n)}
  \qquad \forall \sigma : \Delta_n \to \Delta_1 ,
\]
\[
  \underbrace{\varphi_{n'}\, A_*(\sigma\tau)}_{A_*(\sigma\tau)\,\varphi \ (\text{par hyp.\ sur } \varphi_{n'})}
  \qquad\text{et}\qquad A_*(\tau)A_*(\sigma)\varphi = A_*(\sigma\tau)\,\varphi .
  \qquad \text{OK.}
\]

\emph{Dém.\ de la prop.} Le cas $n = 0$ est trivial ; pour le cas
$n = 1$ il faut considérer les trois cas $\sigma = \mathrm{id}$, qui donne
$\varphi_1 = \varphi$, puis $\sigma$ application constante de valeur $0$
($\alpha_0$) ou $1$ ($\beta_0$), et cela signifie que
$A_*(\alpha_0) = \alpha$, $A_*(\beta_0) = \beta$, et le fait que
$\varphi_1 = \varphi$ y soit compat.\ \uncertain{tient à ce que} l'hyp.\ \ill{}
sur $\varphi$ \uncertain{qu'elle soit} compatible aux augmentations
$\alpha, \beta$. On peut donc supposer $n \geqslant 2$. Considérons les
\add{\uncertain{corresp.}} applications $\sigma : \Delta_n \to \Delta_1$ ;
elles sont surj.\ ou constantes. Considérons les surjectives, il y en a $n$,
soit $\sigma_r$ ($0 \leqslant r \leqslant n-1$) :
\[
  \sigma_r(0) = 0, \ \ldots, \ \sigma_r(r) = 0, \quad
  \sigma_r(r+1) = 1, \ \ldots, \ \sigma_r(n) = 1 ,
\]
donc
\[
  \begin{cases}
    A_*(\sigma_r)(X_0) = X_0 + \cdots + X_r \overset{\mathrm{df}}{=} U_r \\
    A_*(\sigma_r)(X_1) = X_{r+1} + \cdots + X_n
  \end{cases}
\]
On doit donc avoir
\[
  A_*(\sigma_r)(X_0^{(i)}) = U_r^{(i)} .
\]

\emph{Lemme.} Les $U_r^{(i)}$ ($i \geqslant 1$, $0 \leqslant r \leqslant n-1$)
engendrent la $S$-algèbre $A_n^0$, qui est isomorphe en fait à
$S\{U_0, \ldots, U_{n-1}\}$,

\page{197}
\note{suite directe de la page 195.}
\struck{\ill{}} (donc est la $S$-algèbre engendrée par les
$U_{r,i} = U_r^{(i)}$, $0 \leqslant r \leqslant n-1$, $i \geqslant 1$, soumis aux
relations $U_{r,i}U_{r,j} = c_{ij}\,U_{r,i+j}$). Et
\[
  A_n \simeq S\{U_0, \ldots, U_{n-1}\}[dU_0, \ldots, dU_{n-1}] .
\]
C'est immédiat, puisque \ill{} par $X_0, \ldots, X_{n-1}$ se
\uncertain{réécrivent} par $U_0, \ldots, U_{n-1}$ par
\[
  X_0 = U_0, \quad X_1 = U_1 - U_0, \quad \ldots, \quad X_{n-1} = U_{n-1} - U_{n-2}
\]
et que $A_n^0 \simeq S\{X_0, \ldots, X_{n-1}\}$. Ceci montre l'unicité de
$\varphi_n$, et que on doit avoir
\[
  \varphi_n(\underbrace{U_{r,i}}_{U_r^{(i)}})
  = \varphi_n\bigl(\underbrace{\sigma_r^{*}}_{\overset{\mathrm{df}}{=} A_*(\sigma_r)}(X_0^{(i)})\bigr)
  = \sigma_r^{*}\bigl(\varphi(X_0^{(i)})\bigr)
  = \sigma_r^{*}(\xi_i) ,
\]
i.e.
\[
  \varphi_n(U_r^{(i)}) = \sigma_r^{*}(\xi_i)
  = \sum_j a_{ij}\,U_r^{(j)} + \Bigl(\sum_j b_{ij}\,U_r^{(j)}\Bigr) dU_r .
\]
\note{un premier développement, biffé, est remplacé par la parenthèse.}
Il y a donc un hom.\ d'algèbres \add{unique}
$A_n^0 \xrightarrow{\ \varphi_n^0\ } A_n$ tel que
\[
  \varphi_n^0(U_r^{(i)}) = \sigma_r^{*}(\xi_i) \overset{\mathrm{df}}{=} \xi_{r,i} ;
\]
on doit vérifier les relations
\struck{$\sigma_r^{*}(\xi_i)\sigma_r^{*}($}
\[
  \xi_{r,i}\,\xi_{r,j} = c_{ij}\,\xi_{r,i+j} ,
\]
i.e.
\[
  \sigma_r^{*}(\xi_i)\,\sigma_r^{*}(\xi_j) = c_{ij}\,\sigma_r^{*}(\xi_{i+j}) ;
\]
or $\sigma_r^{*}(\xi_i)\sigma_r^{*}(\xi_j) = \sigma_r^{*}(\xi_i\xi_j)
= \sigma_r^{*}(c_{ij}\,\xi_{i+j}) = c_{ij}\,\sigma_r^{*}(\xi_{i+j})$. OK.

\page{199}
\note{suite directe de la page 197.}
Si on a ainsi étendu $\varphi_n^0$ en un hom.\ d'$S$-algèbres
$\varphi_n : A_n \to A_n$, il faut se donner les images des
$dU_0, \ldots, dU_{n-1}$, et vérifier seulement qu'elles sont de carré nul,
\struck{\ill{}} anticommutent \ill{} entre elles, mais commutent aux $\xi_{s,i}$ :
\[
  \varphi_n(dU_r) = d\varphi_n(U_r) = d\sigma_r^{*}(\xi_1) = \sigma_r^{*}(d\xi_1)
  \qquad
  \Bigl(= d\Bigl(\sum_{j \geqslant 0} a_{1j}\,U_r^{(j)}\Bigr)
  = \Bigl(\sum_{j \geqslant 1} a_{1j}\,U_r^{(j-1)}\Bigr) dU_r\Bigr) ,
\]
\[
  d\xi_1 = \sum_{j \geqslant 1} a_{1j}\,X_0^{(j-1)}\,dX_0 .
\]
\note{au-dessous, la même formule récrite puis biffée.}
Ces éléments sont bien de carré nul \struck{\ill{}}, anticommutent entre eux,
mais commutent aux
\[
  \xi_{s,i} = \sigma_s^{*}(X_0^{(i)})
  = \sum_{j \geqslant 0} a_{ij}\,U_s^{(j)} + \Bigl(\sum_{j \geqslant 0} b_{ij}\,U_s^{(j)}\Bigr) dU_s .
\]
\note{une partie de ce passage est encadrée et rayée de deux traits obliques au crayon ; dans le cadre : $A = S\{U_0, \ldots, U_{n-1}\}[dU_0, \ldots, dU_{n-1}]$, et « vrai si $s \neq r$ », « \ill{} $F + \sum_{0 \leqslant r \leqslant n-1} \ill{}\,dU_r$ ». Le cadre renvoie à une note marginale en oblique.}
\marginal{en oblique : et $\sigma_r^{*}(d\xi_1)$ commute à $\sigma_s^{*}(X_0^{(i)}) = \xi_{s,i}$ (si $s \neq r$) ; \ill{} $\sum a_{1j}\,U_r^{(j-1)}$, $\sum b_{ij}\,U_s^{(j)}$ \ill{} ; $2\bigl(\sum_{j \geqslant 1} a_{1j} \ill{}\bigr)\bigl(\sum_{j \geqslant 0} b_{ij} \ill{}\bigr) = 0$ ; $2a_{11} = 1 + \ill{}$ ; \ill{} $2b_{ik} = 0$ \ill{}}

On utilise, pour prouver que $\varphi_n$ (qui précisément est un hom.\ de
$S$-algèbres) commute à la différentielle, le

\emph{Lemme.} Soient $U_0, \ldots, U_m$ des variables, on considère
\struck{\ill{}} la $S$-algèbre diff.\
\[
  A = S\{U_0, \ldots, U_m\}[dU_0, \ldots, dU_m]
\]
(on munit \uncertain{la} structure de p.d., … $dU_r^{(i)} = U_r^{(i-1)}\,dU_r$ …).
Soit $B$ une $S$-alg.\ diff.\ \struck{\ill{}} (\uncertain{quasi-commutative},
\uncertain{de degré} $\leqslant 1$) \uncertain{graduée mod 2}, différentielle.
On s'intéresse aux hom.\ d'algèbres diff.\ $A \xrightarrow{\ \varphi\ } B$.
Pour $A$ tel, on
\note{la phrase se poursuit au-delà de la fin du lot.}

\end{document}
